\section{Discussion and limitations}\label{sec:discussion}

\subsection{Contribution}\label{sec:contribution}

Checking and revision improved technical correctness over one-shot authoring in the evaluated Python tasks. The application gives the model concrete failure evidence, rechecks revised templates and retains a history for inspection. Revision also produced more technical passes than regeneration, although its advantage remains statistically uncertain.

Saved templates generated further drafts without additional model calls. This avoids repeated inference during expansion, but useful variation remains limited by the template's inputs and program structure. The practical benefit depends on how many distinct, suitable questions can be generated before another template is needed.

\subsection{Limitations}\label{sec:limitations}

The main authoring study covered fifteen requests across three model configurations. This small set of Python tasks limits generalisation to other requests and domains. Revision and regeneration had the same candidate allowance, but their token budgets were not matched. In the separate hosted-model study (SoCLaaS), timeouts, empty responses and truncated outputs also affected completion. Workflows requiring several calls have more opportunities to encounter service failures.

Sampled checking can pass without testing every input, so untested variants may still fail or not terminate. Passing technical checks may also leave some user instructions unchecked. Finding a programming feature in reachable code does not show that understanding it is necessary to answer the question. The title's term verifiable refers to explicit checks and recorded evidence and it does not imply formal verification of every generated variant.

\subsection{Educational quality and cost}\label{sec:educational-cost}

Difficulty and relevance results measure agreement with project references rather than actual learner experience. The difficulty questions informed rubric development, while relevance labels came from textbook objectives. These assessments remain aids to human review. Further independent review and learner data are needed to assess educational suitability. The decision-first cascade reduced reasoning-judge calls, but its accuracy advantage over judge-only assessment remains unproven.

Checking and revision required additional inference and time during authoring. Their cost includes unsuccessful attempts, while human review effort was not measured. Reuse avoids further model calls, but duplicates reduce the number of useful questions obtained. The study did not demonstrate overall financial savings.
